\subsection{Enveloping stellar algebra of an involutive Banach algebra}

Lemma 16. --- Let A be an involutive algebra and let p be a semi-norm on A. The following conditions are equivalent:

\begin{enumerate}
\def\labelenumi{(\roman{enumi})}
\item
  We have \(p(xy) \leqslant p(x)p(y)\) , \(p(x^{*}) = p(x)\) and \(p(x)^{2} = p(x^{*}x)\) for all \(x, y \in \mathrm{A}\) ;
\item
  The set R of elements x of A such that \(p(x) = 0\) is a self-adjoint two-sided ideal of A, and the seminorm on A/R induced by p makes A/R an involutive normed algebra whose completion is a C*-algebra;
\item
  There exists a stellar algebra B and a morphism of involutive algebras \(\varphi\) from A to B such that \(p(x) = \|\varphi(x)\|\) for all \(x\in\mathbf{A}\) .
\end{enumerate}

The implications (i) \(\Longrightarrow\) (ii) \(\Longrightarrow\) (iii) \(\Longrightarrow\) (i) are all elementary.

A semi-norm satisfying the conditions of Lemma 16 will be called a stellar semi-norm on the involutive algebra A.

Let A be a normed involutive algebra and let S be the set of stellar seminorms on A. We have \(p(x) \leqslant \|x\|\) for all \(x \in A\) and every \(p \in S\) (I, p.~104, prop. 2). The map \(x \mapsto \|x\|_* = \sup_{p \in S} p(x)\) is a stellar seminorm on A. It is the largest stellar seminorm on A.

Let R be the set of \(x \in A\) such that \(\|x\|_{*} = 0\) . It is a closed two-sided ideal of A. We denote by Stell(A) the completed stellar algebra of A/R for the norm induced by \(x \mapsto \|x\|_{*}\) (Lemma 16, (ii)). The canonical map from A into Stell(A) is continuous, its image is dense in Stell(A), and its kernel is equal to R.

DEFINITION 8. --- The stellar algebra Stell(A) is called the enveloping stellar algebra of the normed involutive algebra A.

If A is commutative, then Stell(A) is commutative; if A is unital, then Stell(A) is unital.

PROPOSITION 20. --- Let A be a normed involutive algebra and let j be the canonical morphism from A into Stell(A). For every stellar algebra B and for every morphism \(\varphi\) of involutive algebras from A to B, there exists a unique morphism \(\varphi'\) of stellar algebras from Stell(A) into B such that \(\varphi = \varphi' \circ j\) .

We denote by \(x \mapsto \|x\|_*\) the norm on Stell(A). Let R be the kernel of \(j\) . The map \(x \mapsto \| \varphi(x) \|\) is a stellar seminorm on A. We therefore have \(\| \varphi(x) \| \leqslant \| x \|_*\) for all \(x \in A\) . The morphism \(\varphi\) therefore defines, by passing to the quotient, a continuous morphism from A/R to B, which extends by continuity to a morphism \(\varphi'\) from Stell(A) to B satisfying \(\varphi = \varphi' \circ j\) . The uniqueness of \(\varphi'\) results from the fact that the image of \(j\) is dense in Stell(A).

COROLLARY. --- Let A be a commutative involutive Banach algebra and let j be the canonical morphism from A into Stell(A). The map \(\mathsf{X}(j)\) is a homeomorphism from \(\mathsf{X}(\mathrm{Stell}(\mathrm{A}))\) onto the subspace \(\mathsf{X}(\mathrm{A})_{h}\) of \(\mathsf{X}(\mathrm{A})\) consisting of the Hermitian characters of A.

The Hermitian characters of A are the morphisms of involutive algebras from A to the stellar algebra C. Proposition 20 therefore implies that \(\mathsf{X}(j)\) is a bijection from \(\mathsf{X}(\mathrm{Stell}(\mathsf{A}))\) onto \(\mathsf{X}(\mathsf{A})_h\) . Since \(\mathsf{X}(j)\) is a homeomorphism onto its image (cf.~I, p.~10), the corollary follows.

We identify \(\mathsf{X}(\mathrm{Stell}(\mathbf{A}))\) with \(\mathsf{X}(\mathbf{A})_h\) by the map \(\mathsf{X}(j)\) . For every \(x\in \mathbf{A}\) , the map \(\mathcal{G}_{\mathrm{Stell(A)}}(j(x))\) is none other than the restriction to \(\mathsf{X}(\mathbf{A})_h\) of \(\mathcal{G}_{\mathbf{A}}(x)\) .

PROPOSITION 21. --- Let A be an involutive Banach algebra and let j be the canonical morphism from A into Stell(A). The radical of A is contained in the kernel of j.

Let \(x\) be an element of the radical of A. Then \(x^{*}x\) belongs to the radical of A, and therefore \(\mathrm{Sp}_{\mathrm{A}}^{\prime}(x^{*}x) = \{0\}\) (I, p.~5, remark 3). Since \(\mathrm{Sp}_{\mathrm{Stell(A)}}^{\prime}(j(x^{*}x)) \subset \mathrm{Sp}_{\mathrm{A}}^{\prime}(x^{*}x)\) , we therefore have \(\mathrm{Sp}_{\mathrm{Stell(A)}}^{\prime}(j(x)^{*}j(x)) = \{0\}\) , whence \(\| j(x)\| ^2 = \| j(x)^{*}j(x)\| = \varrho (j(x)^{*}j(x)) = 0\) (formula (2) of I, p.~104), and therefore \(j(x) = 0\) .
% semantic-fragments: legacy-input-seam
\relax{}
